\section{Transporte del bloque modular al cociente incidencial}
\label{sec:incidence-transport-new}
El bloque central de APP y la incidencia de caras y vértices proporcionan
dos realizaciones de una misma involución. Su relación se expresa mediante
un entrelazador rectangular, antes de efectuar cualquier interpretación
métrica. La congruencia diagonal adyacente
\[
 k^2\equiv(k+1)^2\pmod9
 \quad\Longleftrightarrow\quad 2k+1\equiv0\pmod9
\]
selecciona \(k=4\) en la ventana declarada. El bloque correspondiente es
\[
 B_c=\begin{pmatrix}7&2\\2&7\end{pmatrix}
 =7I_2+2R_c=9P_++5P_-,\qquad
 R_c=\begin{pmatrix}0&1\\1&0\end{pmatrix},\quad
 P_\pm=\frac{I_2\pm R_c}{2}.
\]
La diferencia entre los dos valores propios es cuatro. El transporte
incidencial determina a continuación cómo se conserva esa diferencia
cuando cambian las multiplicidades de los sectores.

Las caras del cubo se indexan por \((i,\epsilon)\), con
\(i\in\{1,2,3\}\) y \(\epsilon\in\{-1,1\}\); sus vértices se
indexan por \(\nu\in\{-1,1\}^3\). Definimos
\[
 M_{(i,\epsilon),\nu}=\mathbf1_{\{\nu_i=\epsilon\}}.
\]
Sean \(R_F\) la oposición de caras y \(R_V\) la inversión de vértices
\(\nu\mapsto-\nu\). En cada entrada se verifica
\[
 \mathbf1_{\{\nu_i=-\epsilon\}}
 =\mathbf1_{\{(-\nu)_i=\epsilon\}},\qquad R_FM=MR_V.
\]
La matriz \(M\) transporta, por tanto, la involución completa.

\begin{theorem}[Descenso incidencial del bloque central]
Sean \(B_F=7I_6+2R_F\) y \(B_V=7I_8+2R_V\). Entonces
\[
 B_FM=MB_V.
\]
Si \(u=(1,1,1,1,1,1)^{\mathsf T}\),
\(P_u=uu^{\mathsf T}/6\) y \(P_-=(I_6-R_F)/2\), el cociente
\[
 Q_\square=\RR^6/\ker(MM^{\mathsf T})
 \simeq\RR u\oplus E_-
\]
tiene dimensión \(1+3\), y el operador inducido es
\[
 \overline B_F=9P_u+5P_-.
\]
\end{theorem}
\begin{proof}
La primera identidad resulta de multiplicar \(R_FM=MR_V\) por dos
y sumar \(7M\). Cada cara contiene cuatro vértices; dos caras adyacentes
comparten dos y dos caras opuestas no comparten ninguno. De esta cuenta se
obtiene, entrada por entrada,
\[
 MM^{\mathsf T}=12P_u+4P_-.
\]
El primer proyector tiene rango uno y el segundo rango tres; sus imágenes
son ortogonales. Los valores propios son \(12,4,0\), con multiplicidades
\(1,3,2\). La oposición actúa como \(+1\) en \(\RR u\) y como
\(-1\) en \(E_-\), y conserva el núcleo de la matriz de Gram. Así,
\(B_F\) desciende al cociente y actúa por nueve y cinco en esos dos
sumandos. Su polinomio característico es
\((\lambda-9)(\lambda-5)^3\).
\end{proof}

La asignación temporal al eje \(\RR u\) fija la forma firmada
\(J_L=P_--P_u\). Sobre el cociente, \(I_Q=P_u+P_-\), de modo que
\[
 J_L=\frac{7I_Q-\overline B_F}{2}.
\]
Esta igualdad relaciona la forma de signatura \((3,1)\), con la
convención temporal indicada, y el operador modular transportado. La
matriz de Gram inicial conserva su positividad; la firma se introduce
en la realización posterior mediante el signo asignado a cada sector.
Las dilataciones nueve y cinco siguen siendo las del bloque: la identidad
afín determina la forma y no convierte al bloque en una isometría.

\subsection{Agregación aritmética, norma reticular y funcional de ruta}
La conservación del entero previo a la reducción se hace visible al sumar
las dos hojas. Las matrices de agregación modular son
\[
 M_\Pi=\begin{pmatrix}4&5&6\\5&4&6\\6&6&9\end{pmatrix},\qquad
 M_\Sigma=\begin{pmatrix}5&6&4\\6&4&5\\4&5&6\end{pmatrix}.
\]
Sus sumas son cincuenta y uno y cuarenta y cinco. Al reconstruir los
nueve elementos de cada bloque, el entero y su acarreo satisfacen
\[
 2025-810=(459-405)+9(174-45)=54+9\cdot129.
\]
El agregado residual y el cociente constituyen dos coordenadas de esta
igualdad. En la construcción reticular posterior, el valor cincuenta y
cuatro aparece como norma de la cámara marcada
\(v_\alpha=4\rho_1+\rho_2+\cdots+\rho_{12}\), para vectores
ortogonales \(\rho_j^2=2\). En la realización material aparece como
evaluación del funcional de ruta \(D_1\) del perfil \((80,54,6)\).
Cada aparición conserva su aplicación de origen: agregación con acarreo,
forma reticular y evaluación orientada de ruta. La coincidencia de sus
valores adquiere contenido estructural mediante esas aplicaciones, cuyos
dominios y pruebas se desarrollan en sus secciones respectivas.
